\section*{Executive Summary}
\addcontentsline{toc}{section}{Executive Summary}

Nexa Environmental Solutions (Nexa) was retained by Meridian Development Corp. (Meridian) to conduct a Phase I Environmental Site Assessment (ESA) of the Vertex Industrial Facility located at 1005 Industrial Parkway, Commerce City (the ``Subject Property''). The assessment was performed in conformance with the scope and limitations of ASTM International Standard Practice E1527-21 \cite{astme1527} and the U.S. Environmental Protection Agency (EPA) All Appropriate Inquiries (AAI) Rule \cite{epa2018}.

\subsection*{Findings and Site Description}
The Subject Property consists of a 4.5-acre parcel developed with a single-story, 45,000-square-foot industrial facility built in 1974. The site is currently operated by Vertex Assemblies for light electronics assembly and warehousing. Historical records indicate the property was agricultural prior to 1958, a machine shop (operated by Apex Machine Shop) from 1958 to 1974, and a metal plating facility (operated by Synapse Plating Corp.) from 1974 to 1998.

This assessment has revealed evidence of two Recognized Environmental Conditions (RECs), one Controlled Recognized Environmental Condition (CREC), and one De Minimis Condition in connection with the Subject Property:

\begin{recbox}{REC 1: Historical Metal Plating Operations}
Historical electroplating and metal finishing operations conducted by Synapse Plating Corp. on the Subject Property between 1974 and 1998 utilized chlorinated solvents (primarily trichloroethylene, or TCE) and heavy metals. Subsurface releases of volatile organic compounds (VOCs) and metals may have occurred in the former plating room and chemical storage areas.
\end{recbox}

\begin{recbox}{REC 2: Potential Vapor Intrusion from Upgradient Dry Cleaner}
An adjoining property to the north (Nexa Cleaners, located at 1001 Industrial Parkway) is listed in the regulatory databases with a documented release of tetrachloroethylene (PCE) impacting groundwater. Given the upgradient hydraulic position relative to the Subject Property, this release represents a migration pathway and potential vapor intrusion risk to the Subject Property.
\end{recbox}

\begin{infobox}{CREC 1: Closed Petroleum Underground Storage Tank (UST)}
A 5,000-gallon diesel fuel UST located on the eastern portion of the Subject Property was removed and clean-closed in 2018. Soil and groundwater sampling conducted at the time of removal indicated minor residual petroleum hydrocarbons below regulatory action levels. A covenant and deed restriction remains in place limiting the installation of water supply wells on the property.
\end{infobox}

\subsection*{De Minimis Conditions}
Minor, surficial oil staining was observed on the concrete slab near the south loading dock. The staining was isolated to a 2-square-foot area and did not appear to penetrate the concrete or impact soil. This staining is classified as a de minimis condition.

\subsection*{Conclusions and Recommendations}
Based on the identified RECs, Nexa recommends conducting a Phase II Environmental Site Assessment (ESA). The Phase II investigation should include soil boring installations, soil vapor sampling, and groundwater monitoring well installation to evaluate potential impacts from historical plating operations and vapor intrusion from the upgradient dry cleaning facility.
\newpage

